\section{One possible downward curvature edge, with immediate recovery}
The earlier proposed monotonicity of source increments after
$\kappa+3$ is FALSE, even arbitrarily far beyond that rank.
The exact legal append gives a different unbounded statement, proved here.
All assertions concern actual normalized SORTED single arms.

Write $L\ge2$ for the final capacity, $F$ for the ordinary polynomial,
$E=1-z+(2z-1)F$, $w_n=\Delta E_n$, and
$D_0=D_1=0$, $D_n=(F_n-F_{n-1})^+$ for $n\ge2$.
Let $\ell$ be the FIRST mode of $D$ and $R$ its last positive rank,
which is the FIRST ordinary mode of $F$. Extend sequences by zero.

\begin{theorem}[Single curvature notch and immediate recovery]
For every actual nonrectangular sorted arm,
\[
 w_n\ge w_{n-1}\qquad(3\le n\le\ell,\ n\ne L+2).        \tag{1}
\]
If $L+3\le\ell$, then
\[
                         w_{L+3}\ge w_{L+1}.            \tag{2}
\]
Consequently every $L+3\le j\le\ell$ dominates ALL earlier
source increments: $w_i\le w_j$ for every $i\le j$.
For any $t\ge L$ and $i<j\le\ell$ with $j-i=t$, one also has
$w_i\le w_j$, including windows whose head lies at the exceptional notch.
\end{theorem}

\paragraph{Established inputs.}
The positive derivative $D$ is unimodal with positive interval
support $[2,R]$. Its rising edges through $\ell$ are strict except
possibly at $L+1$. The source increment is nonnegative through $R$,
is zero at every rank below $\kappa\ge3$, and
\[
             w_i=2D_{i-1}-D_i\qquad(3\le i\le R).      \tag{3}
\]
The already established TC and the complete second-offset lemma imply
\[
                      w_i\ge1\qquad(\kappa+2\le i\le R). \tag{4}
\]
Indeed use the second-offset lemma at $\kappa+2$ and TC at later
indices through $\ell$. At an index after $\ell$ but at most $R$,
unimodality gives $w_i\ge D_{i-1}\ge1$. The first birth is
$w_\kappa=\delta\ge1$. None of these facts assumes (1).

\begin{proof}
Induct on the number of non-one capacities. After one capacity
$\ell=2$, so the assertion and recovery condition have empty ranges.
In an append step use OLD symbols as above and let $t\ge L$ be the
NEW largest capacity. The exact source recurrence gives
\[
 \widehat w_n-\widehat w_{n-1}
     =w_{n-1}-w_{n-t-1}
       +\mathbf1_{n=t+1}-2\mathbf1_{n=t+2}
                                 +\mathbf1_{n=t+3}.     \tag{5}
\]
For $3\le n\le\widehat\ell$, the entire NEW rising prefix agrees
with $P=z^2[t-1]+S_tD$. Its single-flat theorem therefore implies
the STRICT comparison
\[
                         D_{n-1}>D_{n-t-1}.            \tag{6}
\]
At an ordinary strict edge this follows directly from
$\Delta P_n=D_{n-1}-D_{n-t-1}-\mathbf1_{n=t+1}$.
At the only possible flat edge $n=t+1$, (6) still follows because
the subtracted unit makes the entering-minus-leaving difference one.
In particular $j=n-1\le R$.

Put $i=n-t-1$, $j=n-1$. We prove $w_i\le w_j$.
If $i\le2$, this follows from $w_i=0$ and nonnegativity through $R$.
Suppose $i\ge3$. If $j\le\ell$, then $j=i+t\ge L+3$.
By OLD (1), OLD (2), and chaining before and after the sole notch,
every such $j$ dominates every earlier $i$. This use of the induction
hypothesis is exactly the record consequence stated in the theorem.
If instead $j>\ell$, (6) and unimodality force $i<\ell$.
The OLD rising side gives $w_i\le D_i$, and the OLD falling side gives
$w_j\ge D_j$. Hence $w_j\ge D_j>D_i\ge w_i$ by (6).
This proves the required comparison in every case. Equation (5) now
gives NEW (1): its only negative isolated term is at $n=t+2$.

For NEW recovery suppose $t+3\le\widehat\ell$. At that NEW edge
(6) gives $D_{t+2}>D_2=1$, so $t+2\le R$.
Adding (5) at $t+2$ and $t+3$ gives the exact identity
\[
             \widehat w_{t+3}-\widehat w_{t+1}
                              =w_{t+1}+w_{t+2}-1,       \tag{7}
\]
because $w_1=w_2=0$. If $t=\kappa-1$, its first term
$w_{t+1}=w_\kappa\ge1$ pays the unit. If $t\ge\kappa$,
its second term $w_{t+2}\ge1$ by (4). The other term is
nonnegative, since both indices are at most $R$. Thus (7) is
nonnegative, proving NEW (2).

For completeness, the record consequence uses only the newly proved
two assertions: through $L+1$ chain (1); skip to $L+3$ using (2);
then chain (1) again. The isolated rank $L+2$ is also bounded by
$L+3$ by the nonexceptional adjacent edge. Below rank three use zero
and nonnegativity. For a window of length $t\ge L$, either its leaving
rank is at most two or its head is at least $L+3$. This proves the
spaced comparison and closes the induction without circularity.
\end{proof}

\begin{proposition}[An unbounded actual obstruction to removing the notch]
For all integers $T\ge6$ and $\alpha\ge T$, the actual arm
\[
                              X=1^\alpha(2,T,T+1)
\]
has $\kappa=5$, $\delta=2$, $\eta=6$, FIRST positive-derivative mode
$\ell=T+3$ and height $2T-1$. Its source increments satisfy
\[
                    w_{T+2}=2T-2>2T-3=w_{T+3}.         \tag{8}
\]
Thus even strictly after $\kappa+3$, and at arbitrarily large offsets
with FIXED $\kappa=5$, source curvature need not be nondecreasing.
The fall is exactly at finalcapacity$+2$, as (1) permits.
\end{proposition}
\begin{proof}
The first arm $1^\alpha2$ has positive derivative $z^2$.
Appending $T$ gives exactly
\[
             D=z^2+2\sum_{i=3}^{T}z^i+z^{T+1}+z^{T+2}. \tag{9}
\]
The ordinary splice retains this entire polynomial for $\alpha\ge T$;
at the boundary $\alpha=T$ the next ordinary signed derivative is
negative, and for larger $\alpha$ it is zero. Thus its first derivative
mode is three and its last positive rank is $T+2$.
Appending $t=T+1$ gives $P=z^2[T]+S_{T+1}D$.
Its edges from three through $T+1$ strictly rise; at $T+2$ the edge
is zero, at $T+3$ it rises by one, and at $T+4$ it falls by one.
The remaining tail cannot rise. The actual derivative agrees through
$T+3$ and is bounded above by $P$ at every later rank, so its FIRST
mode is $T+3$, of height $D(1)=2T-1$.
The source phases give $(\kappa,\delta,\eta)=(5,2,6)$.
For the two-capacity OLD arm, (9) gives
$E_{T+1}=2T-3$, $E_{T+2}=2T-2$.
The exact source window for its append at $T+2$ adds one and at
$T+3$ subtracts one. It therefore gives (8).
\end{proof}

\begin{corollary}[The complete second overlap when either first weight is three]
When BOTH ordinary factors are strictly rising at the compared edge,
the second source-overlap branch has tail nonrecovery whenever at least
one arm has first missing weight three. The counts, other first weight,
initial runs and capacities are unrestricted. The usual OLD head premise
gives all-head S; the entire zero-head branch for both counts at least two
needs neither C nor OLD A2.
\end{corollary}
\begin{proof}
If a forward ordinary derivative is nonincreasing at the compared edge,
the already proved direct monotonicity branch applies. Suppose both
strictly rise. An actual sorted arm with $\kappa=3$ and first derivative
mode at least three necessarily has $(\delta,\eta)=(1,1)$.
Indeed with one capacity its first derivative mode is two. With two
capacities, first weight three forces both capacities to be two;
its derivative is bounded above by $z^2+z^3+z^4$ and starts at one,
so again its FIRST mode is two. With at least three capacities, first
weight three forces ALL capacities to be two, and the last two stationary
source steps give $(\delta,\eta)=(1,1)$ exactly.

Name this arm $Y$ and the other first weight $p$. At the second overlap
$N=p+4$, the source convolution is exactly
$(E_XE_Y)_N=\delta_X+\eta_X$.
The proved joint source theorem gives
$\Delta E_X(N)\ge\delta_X+\eta_X$ since $N\ge p+4$ and both
positive derivatives strictly rise. The other source increment is
nonnegative. Thus the exact source contrast gives tail halving.
For zero head the previously proved double-defect lower bound forces both
ordinary signed slopes positive when the nonnegative tail is zero or one;
the argument just given or the flat/falling derivative branch therefore
applies without an OLD head premise. The direct zero-head budget is positive.
\end{proof}

The location and immediate recovery of the sole permitted fall are
positive unbounded structural information. General C, second overlaps
with both first weights at least four and higher multiplicities,
arbitrary deeper overlaps, full A2 and the original conjecture remain open.
